\documentclass{article}
\usepackage[utf8]{inputenc}
\usepackage[T1]{fontenc}
\usepackage{geometry}
\geometry{
    a4paper,
    left=2.5cm,
    right=2.5cm,
    top=2.5cm,
    bottom=2.5cm
}
\usepackage{graphicx}
\usepackage{float}
\usepackage{hyperref}
\usepackage[english]{babel}
\usepackage{amsmath}
\usepackage{amssymb}
\usepackage{array}
\usepackage{booktabs}
\usepackage{caption}
\usepackage{subcaption}
\usepackage{enumitem}
\usepackage{fancyhdr}
\usepackage{titlesec}
\usepackage{url}
\usepackage{textcomp}
\usepackage{xcolor}
\usepackage{listings}
\usepackage{multirow}
\usepackage{longtable}
\usepackage{tabularx}
\usepackage{times}

% Set up page style
\pagestyle{fancy}
\fancyhf{}
\rfoot{\thepage}
\renewcommand{\headrulewidth}{0pt}

% Title formatting
\titleformat{\chapter}[display]{\normalfont\bfseries}{\chaptertitle\thechapter}{0pt}{\Huge}
\titlespacing*{\chapter}{0pt}{50pt}{40pt}

% Document info
\title{Template Document}
\author{Academic Support Agent}
\date{\today}

\begin{document}
\maketitle

%  
    \kaishu \zihao{-5}
	
	%  
	\begin{frame}
		\titlepage
		%         logo Logo
		%\begin{figure}[htpb]
		%	\begin{center}
		%		\includegraphics[width=0.2\linewidth]{SCU_Logo.png} 
		%	\end{center}
		%\end{figure}
	\end{frame}
	
	%  
	\begin{frame}
		
	\end{frame}
		
	
	%---------------------------------------------------------
	%  
	\section{ }
	
		\begin{frame}
			\begin{itemize}[<+-| alert@+>] 
	    % structure alert 
				\item  \LaTeX{} Beamer ,  overleaf 
				\item   Xe\LaTeX{}   \qquad  overleaf  Xe\LaTeX{} 
				\item Overleaf  \url{https://www.overleaf.com/latex/templates/} 
			\end{itemize}
		\end{frame}
	
	%-------------------------------------------------------	
	\section{ }
	    \subsection{Beamer }	
	    
		\begin{frame}
			\begin{enumerate}
				\item<1->   \LaTeX{}  
				\item<2->  Tsinghua 
				\item<3->  overleaf beamer  \newline \url{https://cn.overleaf.com/latex/templates/thu-beamer-theme/vwnqmzndvwyb}
				\item<4->  beamer  
				\newline
				\url{https://www.latexstudio.net/archives/11443.html}
			\end{enumerate}	
		\end{frame}
	
	%-----------------------------------------------------
	\section{ }
	    \subsection{ }
	
	\begin{frame}{ Overleaf THU Beamer Theme }
		\begin{itemize}[<+->]
			\item  logo
			\item  
			\item  lstlisting  \\
			\url{https://cn.overleaf.com/latex/templates/thu-beamer-theme/vwnqmzndvwybl}
			\item  
		\end{itemize}
	\end{frame}
		
	\begin{frame}{ \LaTeX{}  }
		\begin{enumerate}[<+->]
			\item  B  \\ \url{https://www.bilibili.com/video/BV15b411j7Au?spm_id_from=333.999.0.0}
			\item   Windows  Linux  CTEX LaTeX2$\epsilon$ 
			\item  \LaTeX{}   
			\item  CSDN
			\item  GitHub
		\end{enumerate}
	\end{frame}
	
	\subsection{ Beamer}
	
	\begin{frame}
    	\begin{itemize}
    	    \item \LaTeX{}  
    	\end{itemize}
        \begin{table}[h]
    		\centering
    		\begin{tabular}{c|c}
    			Microsoft Word & \LaTeX{} \\
    			\hline
    			  &   \\
    			  &   \\
    			  &   \\
    			  &   \\
    			  &   \\
    			  &   \\
    			  &   \\
    			  &   \\
    			  &   \\
    		\end{tabular}
    	\end{table}
    \end{frame}
	
    \begin{frame}{ }
	    \begin{exampleblock}{ } %   * 
    		\begin{equation*}
    		f(x+\xi)=\sum_{n=0}^{N}\frac{1}{n!}\xi^n\frac{d^nf(x)}{dx^n}+R(N,x) 
    		\label{equ-1}
    		\end{equation*}
    	\end{exampleblock}
	
    	\begin{exampleblock}{ \footnote{\tiny {  $\backslash$mathrm\{\}   $\backslash$text\{\}   $clip$  $\mathrm{clip}$  }}}
    		%   &  
    		\begin{align}
    		    \frac{1}{n!}\int_{H_x}\xi^ng_N^p(\xi)d\xi=\delta_{np}=
    		    \begin{cases}
    			1 & n=p\\
    			0 & n\neq p\\
    		    \end{cases}
    	    	\label{equ-2}
    		\end{align}
    	\end{exampleblock}
    \end{frame}
	
	\begin{frame}
		\begin{exampleblock}{ }
			% Taken from Mathmode.tex
			\begin{multline}
			A=\lim_{n\rightarrow\infty}\Delta x\left(a^{2}+\left(a^{2}+2a\Delta x+\left(\Delta x\right)^{2}\right)\right.\label{eq:reset}\\
			+\left(a^{2}+2\cdot2a\Delta x+2^{2}\left(\Delta x\right)^{2}\right)\\
			+\left(a^{2}+2\cdot3a\Delta x+3^{2}\left(\Delta x\right)^{2}\right)\\
			+\ldots\\
			\left.+\left(a^{2}+2\cdot(n-1)a\Delta x+(n-1)^{2}\left(\Delta x\right)^{2}\right)\right)\\
			=\frac{1}{3}\left(b^{3}-a^{3}\right)
			\end{multline}
		\end{exampleblock}
	\end{frame}
	
	\begin{frame}{ }
	    % Taken from SUSTech beamer 
		\begin{columns}[T] % align columns
			\begin{column}<0->{.40\textwidth}
				\begin{figure}[thpb]
					\centering
					\resizebox{1\linewidth}{!}{
						\includegraphics{JSNU.jpg}
					}
					%	\includegraphics[scale=1.0]{figurefile}
					%	 pdf 	
					\caption{\tiny{Jiangsu Normal University}}
					\label{fig:1}
				\end{figure}
			\end{column}%
			\hfill%
			\begin{column}<0->{.65\textwidth}
				\begin{itemize}
					\item<1->  
					\begin{itemize}
						\item<1-> Jiangsu Normal University
					\end{itemize}
					\item<2->  
					\begin{itemize}
						\item<2-> Xuzhou, Jiangsu Province, P.R. China
					\end{itemize}
				\end{itemize}
			\end{column}%
		\end{columns}
	\end{frame}

	\begin{frame}[fragile]{\LaTeX{}  }
		\begin{exampleblock}{ }
			\centering
			\footnotesize
			\begin{tabular}{llll}
				\cmd{chapter} & \cmd{section} & \cmd{subsection} & \cmd{paragraph} \\
				  &   &   &   \\ \hline
				\cmd{centering} & \cmd{emph} & \cmd{verb} & \cmd{url} \\
				  &   &   &   \\ \hline
				\cmd{footnote} & \cmd{item} & \cmd{caption} & \cmd{includegraphics}  \\
				  &   &   &   \\ \hline
				\cmd{tabular} & \cmd{label} & \cmd{cite} & \cmd{ref} \\
				  &   &   &   \\ \hline
			\end{tabular}
		\end{exampleblock}
		\begin{exampleblock}{ }
			\centering
			\footnotesize
			\begin{tabular}{lll}
				\env{table} & \env{figure} & \env{equation}\\
				  &   &   \\\hline
				\env{itemize} & \env{enumerate} & \env{description}\\
				  &   &   \\\hline
			\end{tabular}
		\end{exampleblock}
	\end{frame}
	
	%-------------------------------------------------
	\section{ }	
	
	\begin{frame}<1->
		 \begin{thm}
	        contents...
	        \label{thm-1}
		 \end{thm}
		% : < >,  <1-3>, <2>, \begin{thm}<2-> 2 , .
		\begin{block}{ 1.1}<2->
            This is a theorem.
            \begin{equation}
                a^2 + b^2 = c^2
            	\notag	%  
            	\label{equ-3}
            	\end{equation}
        \end{block}
    	% beamer theorem corollary definition definitions fact example examples proof proof 
    	
		\begin{proof}<3->
            \textit{Trivial.} 
        \end{proof}
		            
        \begin{corollary}<4->
            This is a corollay.
            \begin{equation}
                c^2 = b^2 + a^2
                \label{equ-4}
            \end{equation}
        \end{corollary}
	\end{frame}
	
	\begin{frame}{ }
		\begin{itemize}
			\item   eps, ps, pdf
			\begin{itemize}
				\item METAPOST, pstricks, pgf $\ldots$
				\item Xfig, Dia, Visio, Inkscape $\ldots$
				\item Matlab / Excel   pdf
			\end{itemize}
			\item   png, jpg, tiff $\ldots$
			\begin{itemize}
				\item  
				\item  
			\end{itemize}
		\item  eps jpg 
		\end{itemize}
		
	\end{frame}

	%--------------------------------------------------
	\section{ }	
	
		\begin{frame}
			\begin{itemize}
				\item 2021.12.29   
				\item 2022.01.06   
				\item 2022.03.24-2022.03.30  
				\item 2022.04.14-2022.04-28  
				\item 2022.05.05   
				\item 2022.05.05-2022.05.10  
				\item 2022.05.09-2022.05.11  
				\item 2022.05.12-2022.05.18  
			\end{itemize}
			  \cite{name1}	
		\end{frame}
	
	%----------------------------------------------
	\section{ }
		
		\begin{frame}[allowframebreaks]
			
			
				
			%  
			% \tiny
            
			 Bib Tex 
		\end{frame}

	%-------------------------------------------
\begin{frame}
    \begin{center}
        {\Huge\calligra Thanks!}
    \end{center}
\end{frame}
\end{document}